\chapter{Johnson Bound}
%%%%%%%%%%%%%%%%%%%%%%%%%%%%%%%%%%%%%%%%%%%%%%%%%%%%%%%%%%%%%%%%%%%%%%%%%%%%%%%
\section{Overview}

This chapter proves the \textbf{binary Johnson bound}: for a binary code $C$
of length $n$ with minimum distance $d$, if every codeword has weight at most
$w \le J_2(n,d)$ where
\[
J_2(n,d) = \frac{n}{2} - \sqrt{\frac{n(n-d)}{4}},
\]
then $|C| \le 2n$.

The proof reduces to a \emph{Rankin bound} from geometry: a set of unit
vectors in $\mathbb{R}^n$ with pairwise non-positive inner products has
size at most $2n$.  Codewords are embedded via a shifted $\{\pm1\}$ map so
that the distance and weight constraints force the inner-product condition.

%%%%%%%%%%%%%%%%%%%%%%%%%%%%%%%%%%%%%%%%%%%%%%%%%%%%%%%%%%%%%%%%%%%%%%%%%%%%%%%
\section{Basic definitions}

\begin{definition}[Binary codeword]
\lean{CodingTheory.Johnson.BitVec}
\leanok
A binary codeword of length $n$ is a function $\mathrm{Fin}\,n \to \mathrm{Bool}$.
\end{definition}

\begin{definition}[Hamming weight and distance]
\lean{CodingTheory.Johnson.wt}
\lean{CodingTheory.Johnson.hdist}
\leanok
The Hamming weight $\mathrm{wt}(x) = |\{i : x_i = 1\}|$; the Hamming distance
$\mathrm{hdist}(x,y) = \mathrm{wt}(x \oplus y)$.
\end{definition}

\begin{definition}[Admissible code]
\lean{CodingTheory.Johnson.AdmissibleCode}
\leanok
\uses{CodingTheory.Johnson.wt, CodingTheory.Johnson.hdist}
A code $C \subseteq \{0,1\}^n$ is $(n,d,w)$-admissible if every pair of distinct
codewords has Hamming distance $\ge d$ and every codeword has weight $\le w$.
\end{definition}

\begin{definition}[Johnson radius and shift parameter]
\lean{CodingTheory.Johnson.J2}
\lean{CodingTheory.Johnson.alpha}
\leanok
\[
J_2(n,d) = \frac{n}{2} - \sqrt{\frac{n(n-d)}{4}}, \qquad
\alpha(n,d) = \frac{1}{\sqrt{n(n-d)^{-1}} + 1} - \frac{1}{\sqrt{n/d}+1}.
\]
\end{definition}

%%%%%%%%%%%%%%%%%%%%%%%%%%%%%%%%%%%%%%%%%%%%%%%%%%%%%%%%%%%%%%%%%%%%%%%%%%%%%%%
\section{Embedding and inner-product bounds}

\begin{definition}[Shifted $\pm1$ embedding]
\lean{CodingTheory.Johnson.pmOne}
\lean{CodingTheory.Johnson.shifted}
\leanok
Map $x \in \{0,1\}^n$ to the vector $\hat{x}^\alpha_i = (-1)^{x_i} - \alpha$
(with $\alpha \ge 0$), viewed in $\mathbb{R}^n$.
\end{definition}

\begin{lemma}[Inner product formula for $\pm1$ vectors]
\lean{CodingTheory.Johnson.inner_pmOne_pmOne}
\leanok
\uses{CodingTheory.Johnson.pmOne, CodingTheory.Johnson.hdist}
$\langle\hat{x},\hat{y}\rangle = n - 2\,\mathrm{hdist}(x,y)$.
\end{lemma}

\begin{lemma}[Inner product bound for shifted vectors]
\lean{CodingTheory.Johnson.inner_shifted_le_expr}
\leanok
\uses{CodingTheory.Johnson.shifted, CodingTheory.Johnson.hdist, CodingTheory.Johnson.wt}
For $\alpha \ge 0$, $\mathrm{hdist}(x,y) \ge d$, and $\mathrm{wt}(x),\mathrm{wt}(y) \le w$,
\[
\bigl\langle\hat{x}^\alpha,\hat{y}^\alpha\bigr\rangle
\;\le\;
(n - 2d) + \alpha^2 n + 2\alpha(2w - n).
\]
\end{lemma}

\begin{lemma}[Johnson arithmetic]
\lean{CodingTheory.Johnson.johnson_arith}
\leanok
\uses{CodingTheory.Johnson.J2, CodingTheory.Johnson.alpha}
If $n > 0$, $1 \le d$, $2d \le n$, and $w \le J_2(n,d)$, then
$(n - 2d) + \alpha^2 n + 2\alpha(2w - n) \le 0$ for the choice $\alpha = \alpha(n,d)$.
\end{lemma}

%%%%%%%%%%%%%%%%%%%%%%%%%%%%%%%%%%%%%%%%%%%%%%%%%%%%%%%%%%%%%%%%%%%%%%%%%%%%%%%
\section{Rankin bound}

\begin{lemma}[Rankin reduction]
\lean{CodingTheory.Johnson.rankin_reduction}
\notready
In a finite inner-product space, if a set $S$ of unit vectors has pairwise
non-positive inner products, then one can project away one vector and
reduce to a smaller set in the orthogonal complement with the same property.
\end{lemma}

\begin{theorem}[Rankin bound]
\lean{CodingTheory.Johnson.rankin_bound_general}
\notready
\uses{CodingTheory.Johnson.rankin_reduction}
A set of unit vectors in an $n$-dimensional real inner product space
with pairwise non-positive inner products has size at most $2n$.
\end{theorem}

\begin{corollary}[Rankin bound for finite sets in $\mathbb{R}^n$]
\lean{CodingTheory.Johnson.rankin_finset_bound}
\leanok
\uses{CodingTheory.Johnson.rankin_bound_general}
A finite set of unit vectors in $\mathbb{R}^n$ with pairwise non-positive inner
products has cardinality at most $2n$.
\end{corollary}

%%%%%%%%%%%%%%%%%%%%%%%%%%%%%%%%%%%%%%%%%%%%%%%%%%%%%%%%%%%%%%%%%%%%%%%%%%%%%%%
\section{Johnson bound}

\begin{theorem}[Johnson bound, parametric form]
\lean{CodingTheory.Johnson.binary_johnson_card_bound_parametric}
\leanok
\uses{CodingTheory.Johnson.rankin_finset_bound, CodingTheory.Johnson.inner_shifted_le_expr}
Let $C \subseteq \{0,1\}^n$ be a code with minimum distance $\ge d$ and all
weights $\le w$.  For any $\alpha \ge 0$ with
$(n-2d) + \alpha^2 n + 2\alpha(2w-n) \le 0$,
\[
|C| \;\le\; 2n.
\]
\end{theorem}

\begin{theorem}[Johnson bound]
\lean{CodingTheory.Johnson.binary_johnson_card_bound}
\leanok
\uses{CodingTheory.Johnson.binary_johnson_card_bound_parametric, CodingTheory.Johnson.johnson_arith}
If $n > 0$, $1 \le d$, $2d \le n$, and all codewords of $C$ have weight $\le w \le J_2(n,d)$
and minimum distance $\ge d$, then $|C| \le 2n$.
\end{theorem}

\begin{theorem}[Johnson bound for admissible codes]
\lean{CodingTheory.Johnson.binary_johnson_card_bound_of_admissible}
\leanok
\uses{CodingTheory.Johnson.binary_johnson_card_bound, CodingTheory.Johnson.AdmissibleCode}
Every $(n,d,w)$-admissible code $C$ with $w \le J_2(n,d)$ satisfies $|C| \le 2n$.
\end{theorem}

\begin{definition}[Maximum admissible code size]
\lean{A0}
\leanok
\uses{CodingTheory.Johnson.AdmissibleCode}
$A_0(n,d,w)$ is the maximum size of an $(n,d,w)$-admissible code.
\end{definition}

\begin{theorem}[Johnson bound, radius form]
\lean{binary_johnson_bound_radius}
\leanok
\uses{A0, CodingTheory.Johnson.binary_johnson_card_bound_of_admissible}
If $n > 0$, $1 \le d$, $2d \le n$, and $w \le J_2(n,d)$, then
\[
A_0(n,d,w) \;\le\; 2n.
\]
\end{theorem}
